\section{Game Parts}

\begin{rulebox}{General Rule}
Each \textsc{Cromwell's Victory} game includes:
\begin{tightitems}
  \item One 22-inch by 16-inch game map
  \item One sheet of 100 cardboard playing pieces
  \item One 12-page rules booklet
\end{tightitems}
In addition, one six-sided die is required for play and must be supplied by the players.
\end{rulebox}

\ruleslabel{Cases}

\caseheading{2.1} The mapsheet shows the area over which the battle was fought.

A hex grid has been placed over the terrain on the map to control the movement and positioning of the pieces. Each hexagon (hex) is individually numbered and represents an area 200 yards across. Also printed on the map are a Terrain Key, showing what the various colors and symbols on the map mean, and a Turn Track, described below.

\caseheading{2.2} The Turn Track is used to record the passage of time.

This track consists of 12 numbered boxes, each representing a single game turn consisting of 30 minutes of the historical battle. The Game Turn Marker provided in the game is moved along the track to record the passage of time.

\caseheading{2.3} The Terrain Effects Chart (TEC) indicates both the number of movement points a unit spends to enter different types of terrain and the effect on the attacker's combat strength when attacking into, out of, or across certain terrain features.

\caseheading{2.4} The Combat Table is used to find the outcome of attacks by one player's pieces against the opposing player's pieces.

\caseheading{2.5} The Artillery Table is used to find the outcome of bombardments of one player's pieces by the opposing player's artillery pieces.

\caseheading{2.6} The cardboard playing pieces represent the actual military units, leaders, and artillery pieces that took part in the battle.

There are three basic types of playing pieces:

\begin{tightitems}
  \item Units, representing brigades and regiments of Foot, Light Horse, and Heavy Horse.
  \item Leaders, representing the named officers and their staffs that were actually present at the battle.
  \item Artillery pieces, representing groups of three to five guns.
\end{tightitems}

\caseheading{2.7} The numbers and symbols on the pieces describe the characteristics of each piece.

Units contain six items of information: the unit's name, type, set-up hex, combat strength, morale rating, and movement allowance. Leaders contain five items of information: the leader's name, type symbol, set-up hex, command rating, and movement allowance. Artillery pieces contain only two items of information: a type symbol and set-up hex.

\textbf{Combat Strength} is the relative strength of a unit when attacking and defending, expressed in terms of strength points. Each strength point represents between 75 and 100 men, depending on unit type.

\textbf{Morale Rating} is the relative level of training and motivation of a unit, expressed as a numerical rating. It is used when a unit attempts to rally.

\textbf{Movement Allowance} is the maximum number of clear hexes that a unit can enter in a single movement phase, expressed in movement points. More than one MP can be spent to enter a hex depending on the terrain in the hex and in the hexside through which the hex is entered.

\textbf{Command Rating} is the relative leadership ability of a named leader.

\textbf{Type} is the general category into which a piece fits, and determines the piece's capabilities. Unit types are Foot, Light Horse, and Heavy Horse.

The Manchester Heavy Horse unit contains a dragoon symbol on its front face to indicate that the troops composing that regiment were trained to fight as Dragoons as well as Heavy Horse. This unit is treated in all ways as a Heavy Horse unit; however, when defending in a close or train hex, it is considered a foot unit and uses the appropriate foot line of the Combat Table.

\textbf{Name} is the most common name by which a unit or leader was known at the time of the battle, and is used for purposes of identification.

\textbf{Set-up Hex} is the four-digit identification number of the hex in which the piece is placed at the beginning of the game.

\caseheading{2.8} Color indicates to which force each piece belongs.

There are five separate forces in the game, each of which represents one of the smaller regional or national armies that make up the two opposing armies on the battlefield. Whenever the word \emph{force} is used in the rules, it refers to these smaller components of the armies.

\caseheading{2.9} The terms \emph{enemy} and \emph{friendly} are used to describe the relationship and functions of game pieces to each other.

All of the pieces controlled by one player are friendly to each other. The opposing player's pieces are enemy pieces. Frequently, the phase in which a piece can perform game activity is termed a friendly phase (as in, ``a friendly movement phase''). A phase in which the piece cannot perform game activity while enemy pieces can is termed an enemy phase (as in, ``an enemy movement phase'').
